\section{Predictions and discriminating studies}

\subsection{Falsifiable predictions}

P1: A prespecified, physiologically justified early perturbation changes a defined \ensuremath{\alpha_2\delta}/CaV measure before a network outcome changes. A sufficiently precise null across the nominated window weakens that specific bridge. P2: Selective perturbation or rescue of that measure changes the exposure's network effect. Persistence of the full effect despite successful mediator manipulation favors a bypass explanation.

P3: Measured strain acquisition or function improves prospective prediction of selected developmental dimensions beyond delivery label, covariates, and sampling artifacts. Failure on a held-out cohort rejects the nominated microbial contribution; it does not license switching organisms after seeing the outcome. P4: Constitutional effect modification replicates for predefined markers. A null interaction with adequate precision weakens the named susceptibility model.

P5: Independently estimated RCB predicts differential response to responsive versus internally paced work environments. Failure to replicate, poor reliability, or no incremental value beyond conventional measures weakens the construct. P6: AI with explicit completion constraints improves verified output differently from unconstrained AI, conditional on baseline regulatory measurements. More generated text or more branches alone does not count as improved conversion.

P7: Early microbial/immune variables predict age-appropriate germ aversion after prespecified adjustments. Failure weakens the behavioral-immune extension independently of social-reward pathways. P8: The civilization-level inversion should vary by task, opportunity, and resources. A persistent specialist advantage on depth-intensive tasks is compatible with the hypothesis; an absence of the predicted integrative-profile advantage in its nominated environments is adverse evidence.


\subsection{Minimum discriminating experimental program}

Study A: Mechanistic exposure \ensuremath{\times} mediator experiment. In a specified developing neural preparation, choose one exposure and one \ensuremath{\alpha_2\delta} mechanism before testing. Use a factorial design with exposure/control and selective mediator perturbation/control, supplemented by an orthogonal rescue. Measure viability, differentiation, channel localization, cleavage where relevant, release, and network activity in temporal order. Independent cultures or litters, rather than individual cells alone, define biological replication. Mask analysis, preregister a primary endpoint and smallest effect of interest, and justify sample size from pilot variance. This is the first gate for H1. Organoids cannot model the complete vagal–microbial loop.

Study B: Prospective mother–infant cohort. Recruit before delivery; record indication, labor, antibiotics and timing, maternal strain reservoirs, endocrine measurements with assay validation, feeding, and repeated infant strain/function measures. Predefine a small set of dimensional outcomes and a held-out validation sample. Test confounding-only, parallel-path, and candidate mediation models; use sensitivity analyses for missingness and unmeasured confounding. Accessible peripheral measures are not validated substitutes for brain \ensuremath{\alpha_2\delta} state. A cohort alone cannot close that molecular gap.

Study C: Randomized adult environment crossover. Estimate RCB in independent baseline sessions, then cross task environment with AI assistance and completion constraints. Compare blinded ratings of completed, verifiable artifacts, error burden, retention, task switching, recovery and between-team diversity. A gain in individual quality with convergence across outputs is a distinct result, not a single success score. Counterbalance order and account for learning and carryover. This study directly tests the environmental bridge and remains useful even if PRI fails.

A contingent Study D would replicate a narrowly selected microbial intervention with preregistered mediators and longer follow-up. Randomized assignment estimates the intervention effect; engraftment is post-randomization, so comparing engrafters with non-engrafters does not automatically estimate its causal effect. An animal endocrine \ensuremath{\times} microbial design can test interactions after Study A identifies a plausible molecular target, with litter, maternal-care, and surgical-procedure controls. These are proposed designs, not completed or ethically approved protocols.


\subsection{Evidence ledger}

The accompanying ledger assigns each major connection an ID, claim class, design code, preparation/developmental scope, source, unresolved inference, and falsifier. Figure 1 uses those same IDs. C01 identifies specific constitutional biology; C02 summarizes birth signaling; C03 covers microbial inheritance; C04 and C05 separate mouse social-reward mechanisms from human intervention signals; C06 and C07 distinguish synaptogenesis from release regulation; C08 bounds developmental channel reallocation; C09 marks the unproven perinatal-to-hinge bridge; C10 marks the integrated network-to-regulation inference; C11–C15 address regulatory coupling, developmental reorganization, environmental conversion, disgust, and time-indexed feedback; C16 records null epidemiology; C17 records the VMT trial; C18 isolates the unproven RCCX-to-hinge bridge; C19 preserves pharmacological provenance. C21–C27 add the later-life and counterevidence layer: mixed stress-epigenetic observations, task event-rate effects, bounded camouflaging, group chronotype heterogeneity, selected-cohort autonomic findings, the AI quality–diversity tradeoff, and rare severe CACNA2D1 functional evidence.

For proposed and horizon connections, references provide background plausibility or theory provenance, not evidence for the entire arrow. Removal of all E0 connections should visibly disconnect the complete central cascade. That disconnection is an intended demonstration of the present evidence boundary, not an interface defect.
